% ============================================================
%  Section: Experimental Record of the Design Iterations
%  Insert into Chapter 5 (Reinforcement Learning Formulation),
%  after the Diagnosis section and before Open Problems.
%
%  FIGURE SLOTS: every figure is preceded by a comment of the form
%      % >>> RUN: <path to the images folder>
%  Open that folder, pick the figure you want, and replace the
%  filename in \includegraphics. Captions already state the
%  configuration, so they only need the filename swapped.
% ============================================================

\section{Experimental Record of the Design Iterations}
\label{sec:rl_record}

The preceding sections describe what was tried and why each version was
replaced. This section attaches the measurements. Every number below comes
from a saved training run or from the stored output of the notebook that
produced it.

% ------------------------------------------------------------
\subsection{How runs are identified}
\label{sec:rl_record_naming}
% ------------------------------------------------------------

Training runs are stored in directories whose names encode the configuration.
Three naming conventions were used, corresponding to three eras of the reward
function:

\begin{center}
\begin{tabular}{@{}lll@{}}
\toprule
Era & Directory pattern & Reward \\
\midrule
1 & \texttt{t\{t*\}\_lam\{$\lambda$\}\_beta\{$\beta$\}\_K\{K$_{\max}$\}\_N\{N\}}
  & raw error, $-|e| - \lambda K$ \\
2 & \texttt{t\{t*\}\_errtol\{e\}\_werr\{w\}\_wtime\{w\}\_N\{N\}\_c\{c\}}
  & normalised, \eqref{eq:reward_current} \\
3 & as era 2, with \texttt{\_R\{R\}} appended
  & normalised, with the bubble \\
\bottomrule
\end{tabular}
\end{center}

The names record the numerical constants but not the structural choices, i.e.\
which policy architecture, which action set, and which state definition were
in use. Those were recovered from the saved model checkpoints, which store the
observation space, the action space and the policy class. Since the observation
is a fixed-size vector whose length is determined by the state definition, the
observation length divided by the walker count identifies the state version
uniquely; likewise the action-space size divided by the walker count gives
$|\mathcal{A}|$.

In total the record contains 139 distinct configurations, 667 checkpoints and
approximately 9{,}000 saved figures.

% ------------------------------------------------------------
\subsection{Policy architecture}
\label{sec:rl_record_policy}
% ------------------------------------------------------------

\paragraph{Generation 1, fully shared.}
This version did not learn. Over $1.6\times10^5$ timesteps the reached rate
\emph{fell} from $0.70$ to $0.55$ and the mean episode reward fell from $-0.30$
to $-0.45$, while the episode length grew from $194$ to $263$ steps. The
explained variance of the value function also collapsed, from $0.24$ to $0.08$.
The run was stopped early because every indicator was moving in the wrong
direction.

\paragraph{Generation 2, one walker at a time.}
This version worked. Over $10^6$ timesteps the reached rate rose from $0.04$ to
$0.98$, the mean reward from $-95.8$ to $-2.87$, and the episode length fell
from $385$ to $40.4$ steps. The error at the target improved from
$1.14\times10^{-3}$ to $5.63\times10^{-4}$. It was abandoned not because it
performed badly but because it has no mechanism for choosing which walker to
move, as explained in \cref{sec:rl_policy}.

\paragraph{Generation 3b, flattened joint action.}
Over $1.4\times10^7$ timesteps the reached rate rose from $0.71$ to $0.80$.
The error at the target did not improve; it drifted from $1.52\times10^{-5}$ to
$4.30\times10^{-5}$. This is the version still in use.

% >>> RUN: pick from the saved outputs of the notebooks
%          Archives/5.2 Toy_RL_Shared_policy_old.ipynb   (generation 1)
%          Archives/5.4 Toy RL joint updated without multiDiscrete_old.ipynb  (generation 3b)
\begin{figure}[htbp]
    \centering
    \begin{subfigure}[b]{0.48\textwidth}
        \includegraphics[width=\textwidth]{Figures/RL/PLACEHOLDER_gen1}
        \caption{Generation 1, fully shared policy.}
    \end{subfigure}
    \hfill
    \begin{subfigure}[b]{0.48\textwidth}
        \includegraphics[width=\textwidth]{Figures/RL/PLACEHOLDER_gen2}
        \caption{Generation 2, one walker at a time.}
    \end{subfigure}
    \caption{Training behaviour of the first two policy architectures.
    The shared policy degrades; the one-walker policy reaches the target
    reliably but cannot select which walker to move.}
    \label{fig:rl_gen12}
\end{figure}

\begin{remark}
No checkpoint survives from the MultiDiscrete architecture of
\cref{sec:rl_policy}, and none from the period before action masking was
introduced. Every saved run uses a flattened discrete action space with
masking. The evidence for those two iterations is therefore the stored
notebook output rather than a training run.
\end{remark}

% ------------------------------------------------------------
\subsection{State space}
\label{sec:rl_record_state}
% ------------------------------------------------------------

The state definition passed through five versions, each identifiable from the
observation length recorded in the checkpoints:

\begin{center}
\begin{tabular}{@{}clll@{}}
\toprule
Stage & Observation & Contents & Run directories \\
\midrule
1 & $3N+2$ & absolute coordinates, with $u_i$, with target
  & \texttt{runs1}--\texttt{runs8} \\
2 & $3N$   & with $u_i$, \emph{target removed}
  & \texttt{runs9}, \texttt{runs10}, \texttt{runs13} \\
3 & $2N$   & $u_i$ removed as well
  & \texttt{runs14} \\
4 & $2N+1$ & relative, $t^*$ restored, no $x^*$
  & \texttt{runs14}--\texttt{runs17} \\
5 & $2N+2$ & relative, $t^*$ and $x^*$ (current)
  & \texttt{runs21} onward \\
\bottomrule
\end{tabular}
\end{center}

Stage 2 is the experiment referred to in \cref{sec:rl_state} as omitting the
target from the state. Stage 4 is the intermediate version in which $t^*$ was
restored but $x^*$ had not yet been added. The progression confirms the
account given there: the target was first removed, found to prevent the policy
from conditioning on the horizon, and then reintroduced one coordinate at a
time.

% ------------------------------------------------------------
\subsection{Reward, era 1: the \texorpdfstring{$\lambda$--$\beta$}{lambda-beta} form}
\label{sec:rl_record_reward1}
% ------------------------------------------------------------

The two failures described in \cref{sec:rl_reward} are visible directly in the
training curves.

\paragraph{Failure 1: the accuracy term carried no signal.}
With $\lambda = 0.001$ and no $\beta$ rescaling, training over
$3.1\times10^6$ timesteps drove the reached rate from $0.61$ to $1.00$ and the
mean reward from $-0.615$ to $-0.009$, i.e.\ essentially to its maximum. But
the error at the target did not move at all: it began at
$2.04\times10^{-5}$ and ended at $2.06\times10^{-5}$. Meanwhile the episode
length collapsed from $225$ steps to $8.6$. The agent had optimised the reward
completely, and had done so entirely by minimising the step count. This is the
clearest possible demonstration that the accuracy term was inactive.

\paragraph{Failure 3: the ordering of the two branches broke.}
With $\lambda = 0.005$ and $\beta = 100$, the reached rate fell from $0.70$ to
exactly $0$ over $5\times10^6$ timesteps. The episode length rose from $162$
to $500$, which is the step budget, so every episode was ending in timeout. The
mean reward fell from $-1.12$ to $-3.5$ and the explained variance collapsed
from $0.70$ to $2\times10^{-4}$. With $\beta$ this large, timing out scored
better than reaching, and the policy learned to stop reaching. Notably the
error on the episodes that did reach was \emph{improving}, from
$1.64\times10^{-5}$ to $4.24\times10^{-6}$ --- the agent was being rewarded for
accuracy and punished for arriving.

The same behaviour appeared under a curriculum with $\lambda = 0.05$ and
$\beta = 100$: over $1.2\times10^7$ timesteps the reached rate fell from
$0.79$ to $0$, and the policy entropy collapsed to $3.8\times10^{-4}$, meaning
the policy had become entirely deterministic on a losing strategy.

% >>> RUN: Archives/runs1/  and  Archives/runs8/  and  Archives/runs9/
%          e.g. Archives/runs9/t10_lam0.01_beta10000_K500_N45/seed*/images/
\begin{figure}[htbp]
    \centering
    \begin{subfigure}[b]{0.48\textwidth}
        \includegraphics[width=\textwidth]{Figures/RL/PLACEHOLDER_lam_nobeta}
        \caption{$\lambda = 0.001$, no $\beta$. Reward is maximised, error is
        unchanged, episodes shrink to $8.6$ steps.}
    \end{subfigure}
    \hfill
    \begin{subfigure}[b]{0.48\textwidth}
        \includegraphics[width=\textwidth]{Figures/RL/PLACEHOLDER_beta100}
        \caption{$\lambda = 0.005$, $\beta = 100$. The reached rate falls to
        zero; every episode times out.}
    \end{subfigure}
    \caption{The two failure modes of the unnormalised reward. Neither can be
    fixed by choosing a different constant: the first needs a larger accuracy
    weight, the second needs a smaller one.}
    \label{fig:rl_reward_era1}
\end{figure}

\paragraph{The $\beta$ sweep.}
Because the two failures pull in opposite directions, $\beta$ was searched over
five orders of magnitude, from $10^2$ to $10^7$, across the run sets
\texttt{runs8}, \texttt{runs9}, \texttt{runs10}, \texttt{runs21} and
\texttt{runs22}. No value satisfied both requirements simultaneously at every
horizon, which is what motivated the move to a normalised reward.

% >>> RUN: Archives/runs9/  contains beta = 1e3, 1e4, 1e6, 1e7 at t*=10, N=45
%          Archives/runs8/  contains beta = 1e3, 5e3, 1e4, 1e5 at t*=5, N=45
\begin{figure}[htbp]
    \centering
    \includegraphics[width=0.7\textwidth]{Figures/RL/PLACEHOLDER_beta_sweep}
    \caption{Effect of the accuracy weight $\beta$ at fixed $\lambda = 0.01$.
    Small $\beta$ leaves the accuracy term inactive; large $\beta$ makes
    timing out preferable to reaching.}
    \label{fig:rl_beta_sweep}
\end{figure}

% ------------------------------------------------------------
\subsection{Reward, era 2: normalisation}
\label{sec:rl_record_reward2}
% ------------------------------------------------------------

Normalising both terms into $[0,1]$ removed the ordering problem, because both
branches are then bounded below by $-1$ regardless of the horizon. The run sets
\texttt{runs,NEW22}, \texttt{runs,NEW23}, \texttt{runs722} and
\texttt{runs7-23} use this form, with $w_{\mathrm{err}}$ swept over
$\{0.7, 0.8, 0.9, 1.0\}$, $w_{\mathrm{time}}$ over $\{0.1, 0.2\}$ and
$e_{\mathrm{tol}}$ over $\{10^{-5}, 10^{-3}, 0.25\}$.

% >>> RUN: Archives/runs,NEW23/t50_errtol0.001_werr1_wtime0.1_N50_c0/seed7/images/
%          Archives/runs,NEW22/t50_errtol0.001_werr0.7_wtime0.1_N50_c0.1/seed7/images/
\begin{figure}[htbp]
    \centering
    \includegraphics[width=0.7\textwidth]{Figures/RL/PLACEHOLDER_normalised}
    \caption{A representative run under the normalised reward at
    $t^* = 50\,\Delta t$. Unlike \cref{fig:rl_reward_era1}, the reached rate no
    longer collapses.}
    \label{fig:rl_normalised}
\end{figure}

% ------------------------------------------------------------
\subsection{Action masking}
\label{sec:rl_record_masking}
% ------------------------------------------------------------

After masking was introduced the reached rate rose from $0.50$ to $0.60$ over
$2.5\times10^6$ timesteps. The improvement in the reached rate is modest; the
important change is that training progressed at all, whereas before the agent
could stall indefinitely on a rejected action, as described in
\cref{sec:rl_masking}.

% >>> RUN: "5. Toy RL with intermediate reward inter.ipynb"  (masking added, commit b03b820)
\begin{figure}[htbp]
    \centering
    \includegraphics[width=0.7\textwidth]{Figures/RL/PLACEHOLDER_masking}
    \caption{Training with action masking enabled. The livelock described in
    \cref{sec:rl_masking} no longer occurs.}
    \label{fig:rl_masking}
\end{figure}

% ------------------------------------------------------------
\subsection{Termination: the bubble}
\label{sec:rl_record_bubble}
% ------------------------------------------------------------

Replacing exact arrival with the bubble criterion produced the largest single
improvement in the record. Over $2.1\times10^7$ timesteps the reached rate was
$1.00$ throughout, the error at the target fell from $2.12\times10^{-3}$ to
$6.88\times10^{-4}$, the mean reward improved from $-0.134$ to $-0.063$, and
the episode length halved from $3.65\times10^3$ to $1.79\times10^3$ steps.
Unlike exact arrival, the bubble does not require a walker to land on a single
cell, so the terminal reward is reached in essentially every episode and the
agent receives a usable signal from the start.

% >>> RUN: Taylor linear predictor/runs7-24/t100_errtol0.001_werr1_wtime0.1_N90_c0.0001_R4/seed10/images/
\begin{figure}[htbp]
    \centering
    \begin{subfigure}[b]{0.48\textwidth}
        \includegraphics[width=\textwidth]{Figures/RL/PLACEHOLDER_bubble_early}
        \caption{Early in training.}
    \end{subfigure}
    \hfill
    \begin{subfigure}[b]{0.48\textwidth}
        \includegraphics[width=\textwidth]{Figures/RL/PLACEHOLDER_bubble_late}
        \caption{After convergence.}
    \end{subfigure}
    \caption{Walker trajectories under the bubble termination criterion at
    $t^* = 100\,\Delta t$, $N = 90$, $R = 4$. The dashed box is the bubble
    $\mathcal{B}_R(z^*)$ and the star is the target.}
    \label{fig:rl_bubble}
\end{figure}

% ------------------------------------------------------------
\subsection{Behaviour across horizons}
\label{sec:rl_record_horizon}
% ------------------------------------------------------------

The horizon $t^*$ is the parameter that most affects whether training
converges. Short horizons converge quickly and reliably. Long horizons converge
only in some configurations, and the configurations that succeed are those with
a larger bubble and more walkers.

\begin{table}[htbp]
\centering
\begin{tabular}{@{}rrrrrrrr@{}}
\toprule
$t^*$ & $e_{\mathrm{tol}}$ & $N$ & $c_{\mathrm{shape}}$ & $R$ & steps &
reached & error on reach \\
\midrule
10  & $10^{-3}$ & 30 & $10^{-3}$ & 3 & 4M   & 1.000 & $2.4\times10^{-4}$ \\
10  & $10^{-6}$ & 30 & $10^{-3}$ & 3 & 10M  & 1.000 & $7.0\times10^{-5}$ \\
\midrule
50  & $10^{-3}$ & 30 & $10^{-3}$ & 3 & 19M  & 0.942 & $1.1\times10^{-1}$ \\
50  & $10^{-3}$ & 35 & $10^{-3}$ & 3 & 100M & 0.724 & $6.2\times10^{-3}$ \\
50  & $10^{-5}$ & 50 & $10^{-3}$ & 3 & 100M & 0.998 & $1.5\times10^{-2}$ \\
50  & $10^{-3}$ & 50 & $10^{-2}$ & 4 & 5M   & 0.981 & $1.7\times10^{-1}$ \\
50  & $10^{-3}$ & 50 & $0$       & 4 & 10M  & \textbf{0.999} & $\mathbf{3.6\times10^{-3}}$ \\
\midrule
100 & $10^{-3}$ & 90 & $10^{-4}$ & 4 & 21M  & \textbf{1.000} & $\mathbf{2.2\times10^{-3}}$ \\
\bottomrule
\end{tabular}
\caption{Training outcomes across horizons. Values are averaged over the last
hundred logged points. The two best rows both use $R = 4$; every row with
$R = 3$ is worse, both in reached rate and in error.}
\label{tab:rl_horizons}
\end{table}

Two observations follow. First, the claim that long horizons do not converge is
too strong: at $t^* = 50$ and $t^* = 100$ the policy reaches the target in
essentially every episode, provided the bubble radius is $R = 4$ rather than
$R = 3$ and the walker count is large enough. Second, the configurations that
fail are exactly those in which a walker is most likely to approach the target
alone rather than as part of a group, which is consistent with the stability
argument of \cref{sec:rl_diagnosis}: a larger bubble and more walkers both make
a two-sided stencil easier to obtain.

% >>> RUN: three directories, one per horizon --- pick the "final" figure from each
%   t*=10 : Taylor linear predictor/runs7-27/t10_errtol0.001_werr1_wtime0.1_N30_c0.001_R3/seed6/images/
%   t*=50 : Taylor linear predictor/runs7-24/t50_errtol0.001_werr1_wtime0.1_N50_c0_R4/seed10/images/
%   t*=100: Taylor linear predictor/runs7-24/t100_errtol0.001_werr1_wtime0.1_N90_c0.0001_R4/seed10/images/
\begin{figure}[htbp]
    \centering
    \begin{subfigure}[b]{0.32\textwidth}
        \includegraphics[width=\textwidth]{Figures/RL/PLACEHOLDER_t10}
        \caption{$t^* = 10\,\Delta t$, $N = 30$, $R = 3$.}
    \end{subfigure}
    \hfill
    \begin{subfigure}[b]{0.32\textwidth}
        \includegraphics[width=\textwidth]{Figures/RL/PLACEHOLDER_t50}
        \caption{$t^* = 50\,\Delta t$, $N = 50$, $R = 4$.}
    \end{subfigure}
    \hfill
    \begin{subfigure}[b]{0.32\textwidth}
        \includegraphics[width=\textwidth]{Figures/RL/PLACEHOLDER_t100}
        \caption{$t^* = 100\,\Delta t$, $N = 90$, $R = 4$.}
    \end{subfigure}
    \caption{Converged walker trajectories at three horizons. The number of
    walkers that leave the initial condition grows with the horizon; at
    $t^* = 10$ most walkers never move.}
    \label{fig:rl_horizons}
\end{figure}

% ------------------------------------------------------------
\subsection{Additional ablations present in the record}
\label{sec:rl_record_ablations}
% ------------------------------------------------------------

Three further experiments exist in the archive and are worth recording, since
they test choices asserted elsewhere in this chapter.

\paragraph{Single walker.}
The run set \texttt{runs20} trains with $N = 1$ at $t^* \in \{5,10,20,30\}$.
This is the extreme case of the coordination argument of
\cref{sec:rl_policy}: with one walker there is nothing to coordinate, and the
walker must build its own support alone.

\paragraph{Geometry constraint.}
The directory
\texttt{runs\_adv8/t5\_lam0.001\_beta1000\_K400\_N25,geometry\_constrained}
repeats the adjacent unconstrained run with the two-sided coverage test
enabled, and is the direct test of the admissibility relaxation discussed in
\cref{sec:rl_constraints}.

\paragraph{Joint versus sequential scheduling.}
The run sets \texttt{runs14} and \texttt{runs21} contain jointly sampled
schedules, written as hyphenated horizon lists such as
\texttt{t5-7-10}, \texttt{t5-7-10-15} and \texttt{t5-10-15-20}. The run sets
\texttt{runs29} and \texttt{runs30} contain sequential curricula, written as
\texttt{curriculum\_t2}, \texttt{curriculum\_t4}, \texttt{curriculum\_t5},
\texttt{curriculum\_t6}, \texttt{curriculum\_t10} and
\texttt{curriculum\_t15}. Both forms were therefore tested, and the statement
in \cref{sec:rl_training} that only the joint form was tried should be read as
applying to the current reward only.

% >>> RUN (optional, pick any that are useful):
%   N=1              : Archives/runs20/t10_lam0.01_beta1000_N1/seed8/images/
%   geometry ablation: Archives/runs_adv8/t5_lam0.001_beta1000_K400_N25,geometry_constrained/
%   curriculum ladder: Archives/runs30/curriculum_t5_..._seed1/images/ and curriculum_t15_..._seed1/images/
\begin{figure}[htbp]
    \centering
    \begin{subfigure}[b]{0.48\textwidth}
        \includegraphics[width=\textwidth]{Figures/RL/PLACEHOLDER_curr_t5}
        \caption{Curriculum stage $t^* = 5\,\Delta t$.}
    \end{subfigure}
    \hfill
    \begin{subfigure}[b]{0.48\textwidth}
        \includegraphics[width=\textwidth]{Figures/RL/PLACEHOLDER_curr_t15}
        \caption{Curriculum stage $t^* = 15\,\Delta t$.}
    \end{subfigure}
    \caption{Two stages of the sequential curriculum, $\lambda = 0.001$,
    $\beta = 100$. The policy is trained to convergence at the shorter horizon
    before the longer one is attempted.}
    \label{fig:rl_curriculum}
\end{figure}

% ------------------------------------------------------------
\subsection{Summary}
\label{sec:rl_record_summary}
% ------------------------------------------------------------

\begin{table}[htbp]
\centering
\small
\begin{tabular}{@{}p{3.4cm}p{4.1cm}p{6.4cm}@{}}
\toprule
Iteration & Measured outcome & Conclusion \\
\midrule
Generation 1, shared
  & reached $0.70 \to 0.55$; reward $-0.30 \to -0.45$
  & degrades; stopped after $1.6\times10^5$ steps \\
Generation 2, one walker
  & reached $0.04 \to 0.98$; error $1.1\times10^{-3} \to 5.6\times10^{-4}$
  & works, but cannot select a walker \\
Generation 3b, flattened
  & reached $0.71 \to 0.80$; error $1.5\times10^{-5} \to 4.3\times10^{-5}$
  & retained \\
\midrule
$\lambda = 0.001$, no $\beta$
  & reward $\to -0.009$; error unchanged; episodes $225 \to 8.6$
  & accuracy term inactive \\
$\lambda = 0.005$, $\beta = 100$
  & reached $0.70 \to 0.00$; episodes at the budget
  & timing out beats reaching \\
$\beta$ swept $10^2$--$10^7$
  & no value satisfies both
  & motivates normalisation \\
Normalised reward
  & reached rate stable at $t^* = 50$
  & retained \\
\midrule
Action masking
  & reached $0.50 \to 0.60$
  & removes the livelock \\
Bubble termination
  & reached $1.00$; error $2.1\times10^{-3} \to 6.9\times10^{-4}$
  & largest single improvement \\
\midrule
$R = 3$ vs $R = 4$
  & every $R = 4$ run outperforms every $R = 3$ run
  & bubble radius is decisive at long horizons \\
\bottomrule
\end{tabular}
\caption{Summary of the design iterations and their measured outcomes.}
\label{tab:rl_ledger}
\end{table}
